\documentclass[12pt]{article}
\usepackage[utf8]{inputenc}
\usepackage[T1]{fontenc}
\usepackage[margin=1in]{geometry}
\usepackage{setspace}
\usepackage{xcolor}
\usepackage[colorlinks=true, urlcolor=blue!60!black, linkcolor=blue!60!black]{hyperref}
\setlength{\parindent}{2em}
\setlength{\parskip}{1.25em}
\renewcommand{\baselinestretch}{1.0}
\begin{document}
\begin{center}
{\LARGE \textit{Moore v. City of East Cleveland}, 431 U.S. 494 (1977)\par}
\vspace{0.8em}
{\large Case Brief --- Four-Page Edition\par}
\end{center}
\vspace{1em}

\section{Memory Jogger}

In a fractured 5--4 decision, the Court reversed a grandmother's criminal conviction for living with two grandsons who were first cousins; a four-Justice plurality held that substantive due process protects extended-family living arrangements as deeply rooted in the Nation's history and tradition (\href{https://supreme.justia.com/cases/federal/us/431/494/}{Justia}).

\section{Detailed Case Facts}

East Cleveland's Housing Code limited each dwelling unit to a single ``family,'' defined narrowly as the head of household, a spouse, their unmarried children, their parents, and at most one dependent child with that child's own spouse and children. A separate provision capped occupancy by floor area. Inez Moore, a 63-year-old Black grandmother, lived in her own house with her son, Dale, Sr., and two grandsons, Dale, Jr., and John, Jr., who were cousins; John had come to live with her after his mother died when he was an infant. In 1973 the city told her to remove 10-year-old John as an ``illegal occupant.'' She did not seek a hardship variance and refused to expel him, and the city filed criminal charges.

\section{Procedural History}

The municipal court rejected her constitutional challenge, convicted her, and sentenced her to five days in jail and a \$25 fine. The Ohio Court of Appeals affirmed in reliance on \textit{Village of Belle Terre v. Boraas}, and the Supreme Court of Ohio dismissed her appeal. The U.S. Supreme Court reversed on May 31, 1977.

\section{Judicial Votes}

Reversed, \textbf{5--4}, with no majority opinion.
\begin{itemize}
\item \textbf{Plurality:} Justice Powell, joined by Justices Brennan, Marshall, and Blackmun.
\item \textbf{Concurring:} Justice Brennan, joined by Justice Marshall.
\item \textbf{Concurring in the judgment:} Justice Stevens.
\item \textbf{Dissenting:} Chief Justice Burger; Justice Stewart, joined by Justice Rehnquist; Justice White.
\end{itemize}

\section{Holding}

An ordinance criminalizing cohabitation among close blood relatives violates the Fourteenth Amendment. For the plurality, when government slices into the family itself, courts must examine the importance of the state's interests and how well the regulation serves them, because the extended family is deeply rooted in history and tradition. Justice Stevens's narrower fifth vote held that the unprecedented ordinance bore no substantial relation to health, safety, morals, or welfare under \textit{Euclid} and took Mrs. Moore's property without due process. Powell's history-and-tradition test was later adopted by majorities in \textit{Glucksberg} (1997) and \textit{Dobbs} (2022).

\section{Analysis of Opinions}

\textbf{Plurality (Powell).} \textit{Belle Terre} upheld a limit on unrelated occupants but let any relatives live together; this ordinance intruded on the family itself, so deference was inappropriate. Its relation to overcrowding, traffic, and school costs was ``tenuous'': it allowed a household with a dozen children of one son but barred two cousins, and floor-area limits already addressed density. Acknowledging that substantive due process has been a treacherous field, Powell drew limits from history and basic values, not from an arbitrary line at the nuclear family. A city may not standardize its residents by forcing all into narrowly defined family patterns. Footnote 5 rejected the exhaustion objection.

\textbf{Brennan (concurring).} Confining protection to the nuclear family reflects ``cultural myopia''; extended households are vital to working-class, immigrant, and especially Black families.

\textbf{Stevens (concurring in the judgment).} A homeowner may use her property as she sees fit; no zoning precedent had sustained barring an owner's own relatives.

\textbf{Burger (dissenting).} Mrs. Moore's refusal to seek a variance should bar her claims.

\textbf{Stewart (dissenting).} \textit{Belle Terre} controlled; the family cases protect marrying and raising one's own children, not sharing a kitchen with extended relatives. East Cleveland was itself a predominantly Black city.

\textbf{White (dissenting).} Substantive due process lacks textual and historical roots, and living with more than one set of grandchildren is not implicit in ordered liberty.

\section{Examples: Future Applications}

\begin{enumerate}
\item \textbf{Same-Side: Aunt and Sister.} A city prosecutes a woman for housing her elderly aunt and widowed sister under a parents-and-children definition of family. The law slices into close kinship.
\item \textbf{Same-Side: Grandchildren of Two Children.} A code allows four grandchildren from one married child but not two from different unmarried children---an arbitrary line like East Cleveland's.
\item \textbf{Opposite-Side: College Roommates.} A ban on more than three unrelated occupants, leaving relatives untouched, is governed by \textit{Belle Terre}.
\item \textbf{Opposite-Side: Floor-Area Limits.} A family-blind requirement of 150 square feet per occupant is a neutral density rule both Powell and Stevens endorsed.
\item \textbf{Fence-Sitter: A Functional Family.} An unmarried couple, a partner's sibling, and an unadopted godchild raised for ten years. Courts split on whether chosen families get \textit{Moore}'s protection.
\end{enumerate}

\section{Critique}

\textbf{Originalist views.} Critics echo White that due process guarantees procedure, and note that the history-and-tradition test still lets judges choose the level of generality. Yet later originalists embraced the same test in \textit{Glucksberg} and \textit{Dobbs} as a limit on unenumerated rights, and property scholars favor Stevens's approach.

\textbf{Progressive views.} Scholars praise \textit{Moore} for recognizing the constitutional dignity of multigenerational and kinship-care households, but family-law critics see a Trojan horse: tying protection to blood and marriage entrenched a definition of family later used against unmarried partners and chosen households, and the test became the engine for curtailing autonomy rights.

\section{Key Quotations}

\begin{quote}
``Our decisions establish that the Constitution protects the sanctity of the family precisely because the institution of the family is deeply rooted in this Nation's history and tradition.''
\end{quote}
\noindent --- Justice Powell, plurality opinion, 431 U.S. at 503.

\begin{quote}
``The Judiciary, including this Court, is the most vulnerable and comes nearest to illegitimacy when it deals with judge-made constitutional law having little or no cognizable roots in the language or even the design of the Constitution.''
\end{quote}
\noindent --- Justice White, dissenting, 431 U.S. at 544.

\begin{quote}
``East Cleveland, in contrast, has chosen to regulate the occupancy of its housing by slicing deeply into the family itself.''
\end{quote}
\noindent --- Justice Powell, plurality opinion, 431 U.S. at 498.

\end{document}
